\documentclass[12pt]{article}
\usepackage{amsmath, amssymb, geometry, graphicx}
\geometry{margin=1in}
\title{A Hypothesis for a Sub-Planckian Timescale Based on the Holographic Entropy of the Cosmos}
\author{Naser Ahani \\ Department of Physics, \\ Amirkabir University of Technology (Tehran Polytechnic), \\ Tehran, Iran \\ M.Sc. Student in Fundamental Particles and Field Theory \\ Email: naser.ahani@aut.ac.ir}
\date{November 2025}

\begin{document}
\maketitle

\begin{abstract}
This paper proposes a hypothesis that physical reality operates on a deterministic timescale smaller than the Planck time, defined as the ``Planck-Pulse scale'' ($T_{pp}$). The hypothesis links this sub-Planckian resolution to the total holographic informational degrees of freedom of the universe ($N_Z$). Using the Holographic Principle and the cosmological constant ($\Lambda$) as the defining boundary of the observable universe, we derive a time-independent formula for $N_Z$, yielding $N_Z \approx 3.33 \times 10^{122}$. This leads to a fundamental timescale $T_{pp} \approx 1.62 \times 10^{-166}\,s$. The framework offers a potential bridge between deterministic and probabilistic descriptions of nature, suggesting that quantum uncertainty may emerge from informational coarse-graining of sub-Planckian processes.
\end{abstract}

\section{Introduction: The Problem of Unification and a Sub-Planckian Hypothesis}
Modern physics is structured upon two pillars: General Relativity (GR) and Quantum Mechanics (QM). Each theory has been extraordinarily successful in its own domain, yet their conceptual foundations remain incompatible. Singularities such as the Big Bang and black holes expose this incompatibility most clearly. Numerous unification attempts---string theory, loop quantum gravity, and emergent spacetime models---have made progress but remain incomplete.

Here, we hypothesize that the GR/QM duality is not fundamental, but statistical: both arise from a deeper deterministic substrate of reality that evolves at an ultrafast, sub-Planckian scale. This substrate is characterized by a fundamental timescale $T_{pp}$ defined as
\begin{equation}
T_{pp} = \frac{T_p}{N_Z},
\end{equation}
where $T_p = \sqrt{\hbar G / c^5}$ is the Planck time, and $N_Z$ is the total number of informational degrees of freedom of the cosmos. Our goal is to determine $N_Z$ from known physical principles.

\section{Methodology: The Holographic Principle}
According to the Holographic Principle, all information within a volume can be represented by degrees of freedom on its boundary. The entropy of a black hole, as shown by Bekenstein and Hawking, is proportional to its event horizon area:
\begin{equation}
S_{BH} = \frac{k_B A}{4 l_p^2}, \quad l_p^2 = \frac{\hbar G}{c^3}.
\end{equation}
This implies that information is fundamentally a two-dimensional property encoded on surfaces rather than volumes.

If this principle extends to the entire universe, the total information capacity---and hence $N_Z$---is determined by the area of the cosmic (de Sitter) horizon.

\section{Derivation of $N_Z$ from the Cosmological Constant}
The cosmological constant $\Lambda$ defines a de Sitter horizon with radius
\begin{equation}
R_H = \sqrt{\frac{3}{\Lambda}},
\end{equation}
so that the horizon area is
\begin{equation}
A_H = 4\pi R_H^2 = \frac{12\pi}{\Lambda}.
\end{equation}
Each fundamental unit of information corresponds to a Planck area $4l_p^2 = 4\hbar G / c^3$. Therefore,
\begin{equation}
N_Z = \frac{A_H}{4l_p^2} = \frac{3\pi c^3}{\Lambda \hbar G}.
\end{equation}
Using cosmological parameters $\Lambda \approx 1.088 \times 10^{-52}\, m^{-2}$, we obtain $N_Z \approx 3.33 \times 10^{122}$.

\section{Results and Interpretation}
Substituting this into Eq.~(1) with $T_p \approx 5.39 \times 10^{-44}\,s$, we find
\begin{equation}
T_{pp} = \frac{T_p}{N_Z} \approx 1.62 \times 10^{-166}\,s.
\end{equation}
This timescale suggests that reality, if deterministic, operates through discrete causal ``ticks'' vastly below the Planck limit. Quantum indeterminacy may then reflect the inability of current theories to resolve events occurring on these ultrafast scales.

\section{Discussion and Implications}
The existence of a deterministic sub-Planckian timescale challenges the conventional boundary of physical law. Rather than viewing the Planck time as an absolute limit, we may interpret it as a statistical horizon---a coarse-grained boundary beyond which the fine structure of reality is hidden.

Such a view implies that the apparent randomness of quantum mechanics could arise from informational averaging over a vast number ($N_Z$) of underlying deterministic processes. Time itself, in this interpretation, becomes an emergent phenomenon---a manifestation of discrete causal updates occurring at the Planck-Pulse scale. This perspective aligns with philosophical approaches that treat temporality not as a continuous external parameter but as an intrinsic property of change, intimately tied to information and causation.

While speculative, this framework provides a conceptual bridge between physics and ontology: it situates physical law within a broader informational structure, suggesting that spacetime and matter emerge from deeper patterns of causal organization.

\section{Future Work: Toward a Deterministic Foundation of Spacetime}
Future studies could explore whether the statistical behavior of quantum mechanics can be reconstructed from ensembles of sub-Planckian deterministic events. Developing a thermodynamics or statistical mechanics of spacetime, based on the total number of microstates $N_Z$, could illuminate how classical spacetime geometry and quantum fields emerge as effective descriptions.

A natural extension is to investigate models where information flow or oscillatory dynamics constitute the fundamental substrate of existence. Such a theory might redefine the role of time, treating it as both the medium and the measure of physical processes. This philosophical trajectory, while grounded in physics, points toward a unified view of reality in which causality, information, and existence form an inseparable triad.

\section*{References}
\begin{enumerate}
\item Penrose, R. (2004). \textit{The Road to Reality: A Complete Guide to the Laws of the Universe}. Knopf.
\item Green, M. B., Schwarz, J. H., \& Witten, E. (1987). \textit{Superstring Theory}. Cambridge University Press.
\item Susskind, L. (2005). \textit{The Cosmic Landscape}. Little, Brown and Company.
\item 't Hooft, G. (1993). Dimensional Reduction in Quantum Gravity. \textit{arXiv:gr-qc/9310026}.
\item Susskind, L. (1995). The World as a Hologram. \textit{Journal of Mathematical Physics}, 36(11), 6377--6396.
\item Bekenstein, J. D. (1973). Black Holes and Entropy. \textit{Phys. Rev. D}, 7(8), 2333--2346.
\item Planck Collaboration (2020). Planck 2018 results. VI. Cosmological parameters. \textit{A\&A}, 641, A6.
\item Carroll, S. M. (2001). The Cosmological Constant. \textit{Living Reviews in Relativity}, 4(1), 1.
\item Egan, C. A., \& Lineweaver, C. H. (2010). A Larger Estimate of the Entropy of the Universe. \textit{Astrophys. J.}, 710(2), 1825--1834.
\item Henson, J. (2009). The causal set approach to quantum gravity. In \textit{Approaches to Quantum Gravity}. Cambridge University Press.
\item Verlinde, E. (2011). On the Origin of Gravity and the Laws of Newton. \textit{JHEP}, 2011(4), 29.
\end{enumerate}

\end{document}
